\documentclass{article}
\usepackage[utf8]{inputenc}
\usepackage[T1]{fontenc}
\usepackage{geometry}
\usepackage{graphicx}
\usepackage{float}
\usepackage{hyperref}
\usepackage[english,indonesian]{babel}

\title{Optimasi Pemrosesan Template LaTeX Menggunakan Pendekatan Berbasis AI}
\author{Academic Support Agent}
\date{\today}

\begin{document}
\maketitle

\begin{abstract}

\end{abstract}

\chapter{PENDAHULUAN}


\chapter{TINJAUAN PUSTAKA}


% Original content (commented out):
% \documentclass[oneside,a4paper,11pt]{article}
 \usepackage[brazil,portuges]{babel}
\usepackage[utf8]{inputenc}
\usepackage[T1]{fontenc}
\usepackage{helvet}
\usepackage{verbatim}
\usepackage{xcolor}
\usepackage{pdfpages}
\renewcommand{\familydefault}{\sfdefault}%fonte arial
\usepackage{amsmath, amsfonts, amssymb}
\usepackage[lmargin=2cm,tmargin=2cm,bmargin=2cm,rmargin=2cm]{geometry}%espaços da página
\usepackage{wasysym} %caixas de seleção

\usepackage{etoolbox}
\patchcmd{\thebibliography}{*}{}{}{}

\newcommand{\orientacao}[1]{{\textcolor{blue}{#1}}}

\begin{document}

\thispagestyle{empty}
\begin{center}
\textbf{\orientacao{Coloque aqui o título da proposta }}
\vspace{.5cm}

\end{center}   

\noindent Proponente:

\noindent\textbf{Coloque aqui seu nome}

\vspace{0.5cm}

\noindent Orientador:

\noindent\textbf{Prof. Dr./Me. Nome completo do orientador}

\vspace{0.5cm}

\noindent Coorientador ou Supervisor: \orientacao{deixe apenas a opção que for adequada ao seu trabalho. Em caso de dúvidas, consulte o Art. $8^o$ do Regulamento Complementar das Atividades de Projeto Final de Curso.} \normalsize

\noindent\textbf{Prof. Dr./Me. Nome completo do coorientador}. 

\orientacao{Em caso de supervisão consultar a tabela de títulos CONFEA/CREA. Caso o trabalho não tenha coorientação/supervisão, deixe o espaço em branco.}

\vspace{0.5cm}

\noindent Escreva aqui o resumo da proposta de trabalho (200 a 300 palavras). 

\vspace{.5cm}

\noindent\textbf{Palavras-chave:} \orientacao{coloque aqui cinco palavras-chave que descreva seu trabalho.}

\begin{table}[!h]
\begin{minipage}{1\textwidth}
\textbf{Indicação da(s) grande(s) área(s)}:
  \begin{itemize}
  \item[{\LARGE $\Box$}] Eletrônica;
  \item[{\LARGE $\Box$}] Computação;
  \item[{\LARGE $\XBox$}] Controle;
  \item[{\LARGE $\Box$}] Mecânica.
  \end{itemize}

\vspace{2cm}

\begin{center}
Divinópolis, \today .
\end{center}
\end{minipage}
\end{table}

\newpage

% Corpo do projeto proposto, no máximo, 4 (quatro) páginas distribuídas em: (i) apresentação do problema e dos objetivos do trabalho; (ii) indicação da metodologia de trabalho; (iii) resultados e impactos esperados; (iv) recursos e infraestrutura necessários; e (v) referências bibliográficas.

\section{Apresentação do Problema}\label{Apresentação do Problema}

\orientacao{Nesta seção deverá ser introduzido o problema que será abordado/tratado no projeto final de curso. Lembre-se que nessa seção faz-se necessário a realização de uma boa revisão bibliográfica. Portanto, lembre-se de usar referências atuais e que tenham respaldo ao tema abordado trabalho. Para citá-las corretamente ao longo do texto, use o comando} \verb!\cite{}! \orientacao{e o resultado obtido será: \cite{Rosa2017}. Veja na Seção $5$ como ficará a descrição das referências bibliográficas.}

\subsection{Objetivos da Pesquisa}\label{Objetivos da Pesquisa}

\orientacao{Primeiramente, apresente o objetivo geral do trabalho. Na sequência apresente os objetivos específicos, fazendo uso da estrutura descritiva por tópicos conforme apresentado a seguir:}

\begin{itemize}
\item \orientacao{objetivo específico 01;}
\item \orientacao{objetivo específico 02;}
\item \orientacao{objetivo específico 03;}
\item \orientacao{objetivo específico 04.}
\end{itemize}

\section{Metodologia de Trabalho}\label{Metodologia de Trabalho}

\orientacao{Escreva em detalhes a metodologia para desenvolvimento do trabalho. Lembrem-se, nesta seção deverá ser descrito o passo-a-passo sobre como o trabalho será desenvolvido.}

\section{Resultados e Impactos Esperados}\label{Resultados e Impactos Esperados}

Ao final do trabalho espera-se obter: \orientacao{complete com as informações necessárias de acordo com o objetivo do PFC.}

\section{Recursos e Infraestura Necessários}\label{Recursos Necessários}

\orientacao{Descreva em qual(is) espaço(s) o trabalho será desenvolvido. Além disso, apresente quais recursos (materiais, equipamentos e mão de obra terceirizada) serão necessários para desenvolvimento do projeto proposto. Lembre-se, de sempre apontar e certificar a fonte dos recursos financeiros.}

\bibliographystyle{ieeetr}
\bibliography{referencias}

\newpage   
\begin{center}
\textbf{Plano de Trabalho} 
\end{center}

\noindent\textbf{1. Objetivos das Atividades:} 
\orientacao{Coloque aqui o objetivo geral.}

\vspace{.5cm}

\noindent\textbf{2. Descrição das Atividades:} \orientacao{Veja a seguir um exemplo de como fazer a descrição das atividades. Além disso, note a relação entre a descrição e o crongrama para desenvolvimento das atividades -- Tabela \ref{tab:1}.}

\begin{itemize}
\item[1.1] \orientacao{Revisão da Literatura: revisar as técnicas e métodos presente no controlador RGFL.}
\item[1.2] \orientacao{Programação do Controlador: programar em python o controlador adaptativo.}
\item[1.3] \orientacao{Simulação do Controlador: reproduzir exemplos da literatura.}
\item[2.1] \orientacao{Programação do Observador: programar o observador de estados.}
\item[3.1] \orientacao{Formulação do Problema: formular matematicamente a topologia de controle.}
\item[3.2] \orientacao{Fusão das Topologias: programar a nova topologia de controle.}
\item[3.3] \orientacao{Simulação do Novo Controlador: simular e avaliar o desempenho do controlador proposto.}
\item[4.1] \orientacao{Teste Experimental: avaliar experimentalmente a topologia de controle.}
\item[5.1] \orientacao{Divulgação dos Resultados: divulgação dos resultados em congressos e eventos da área.}
\end{itemize}

\vspace{.5cm}
\noindent\textbf{3. Cronograma de Atividades:} \orientacao{lembre-se de verificar o início do semestre letivo. E não se esqueçam que as atividades de PFC devem ser concluídas em até $2$ semestres letivos, que normalmente finalizam-se em dezembro do respectivo ano letivo.}

\begin{table}[!ht]
\begin{center}
{\footnotesize{\begin{tabular}{c|c|c|c|c|c|c|c|c|c|c|c|c} \hline \hline
Atividade ($\downarrow$) Mês.($\rightarrow$) & $1^{\mbox{\scriptsize o}}$  & $2^{\mbox{\scriptsize o}}$  &$3^{\mbox{\scriptsize o}}$
&$4^{\mbox{\scriptsize o}}$ &$5^{\mbox{\scriptsize o}}$ &$6^{\mbox{\scriptsize o}}$ &$7^{\mbox{\scriptsize o}}$
&$8^{\mbox{\scriptsize o}}$ &$9^{\mbox{\scriptsize o}}$ &$10^{\mbox{\scriptsize o}}$ &$11^{\mbox{\scriptsize o}}$
&$12^{\mbox{\scriptsize o}}$\\ \hline
$1.1$ & $\surd$ & $\surd$ & & & & & & & & & & \\ \hline
$1.2$ & &$\surd$& $\surd$ & & & & & & & & & \\ \hline
$1.3$ & & & $\surd$ & $\surd$ & & & & & & & & \\ \hline
$2.1$ & & & & $\surd$ & $\surd$ & & & & & & & \\ \hline
$3.1$ &  &  & & & $\surd$ &$\surd$ & $\surd$ & & & & & \\ \hline
$3.2$ &  &  & & & & &$\surd$&$\surd$ & & & & \\ \hline
$3.3$ & & & & & & & & $\surd$ & $\surd$ & & & \\ \hline
$4.1$ & & & & & & & & & $\surd$& $\surd$ & $\surd$ & \\ \hline
$5.1$ & & & & & & & & & & & $\surd$ & $\surd$ \\ \hline\hline
\end{tabular}}}
\caption{Cronograma de atividades. \label{tab:1}}
\end{center}
\end{table}

\vspace{.5cm}
\noindent\textbf{4. Local de Desenvolvimento das Atividades:}
\orientacao{Identificação do(s) espaço(s) necessário(s) para desenvolvimento do trabalho.}

\vspace{.5cm}
\noindent\textbf{5. Metodologia de Acompanhamento:}
\orientacao{Descreva a metodologia de acompanhamento e frequência de reuniões com o orientador e coorientador. Façam a leitura do Art. $34^o$ do Regulamento Complementar das Atividades de Projeto Final de Curso a fim de evitar o cancelamento da matrícula na disciplina.}

\newpage  
%instruções: Edite e complete as informações no arquivo editável. Solicite a assinatura do professor orientador, digitalize o arquivo, e altere o nome do arquivo na linha de comando a seguir. Lembre-se que o arquivo digitalizado deverá estar na mesma pasta que o arquivo.tex. 
\includepdf{Termo_Compromisso.pdf}
\end{document}

\end{document}